% Results (hard-wrapped at about 100 columns).
\section{Results}
\label{sec:results}

On the displaced test particles, the generic-trained tracker finds 95.52\% of the targets. The
signal-trained tracker finds 93.03\% of the same targets. The difference is 2.49~$\pm$~0.24 points,
in favour of the tracker that never saw the signal\footnote{The quadrature sum assumes that the two
components are independent, which holds because they come from different sources of randomness.}.
The test set holds 13,563 target particles, so the difference is far more than a handful of
particles. Table~\ref{tab:headline} lists the efficiencies together with the fake-track and
duplicate rates.

\begin{table}[tbp]
  \caption{Efficiency, fake-track rate and duplicate rate of the two trackers on displaced targets.}
  \label{tab:headline}
  \centering
  \begin{tabular}{lcc}
    \hline\hline
     & Generic-trained & Signal-trained \\
    \hline
    Efficiency (\%) & 95.52 & 93.03 \\
    Fake-track rate (\%) & 0.48 & 1.32 \\
    Duplicate rate (\%) & 1.3 & 2.6 \\
    \hline\hline
  \end{tabular}
\end{table}

The generic-trained tracker also produces fewer fake tracks: 0.48\% against 1.32\%. Its duplicate
rate is 1.3\%, compared with 2.6\% for the signal-trained tracker. A tracker that finds more
particles usually pays with more fake tracks, so the lead on both metrics is not a trade-off.
Figure~\ref{fig:fake} shows the fake-track rate as a function of the track length. The difference in
fake tracks comes mostly from short tracks of five or six hits.

\begin{figure}[tbp]
  \centering
  \fbox{\parbox[c][0.45\linewidth][c]{0.8\linewidth}{\centering placeholder artwork: fake}}
  \caption{Fake-track rate versus the number of hits on the track. Most fake tracks of the signal-trained tracker are short.}
  \label{fig:fake}
\end{figure}

\subsection{Displacement, species and pile-up}
\label{sec:dependence}

Figure~\ref{fig:eff_disp} shows the efficiency as a function of the transverse displacement of the
production vertex. Close to the beam line, between zero and 2~mm, the two trackers find 98.87\% and
98.28\% of the targets. In the bin from 2 to 5~mm, the efficiencies are 98.37\% and 97.16\%. For
production radii from 5 to 20~mm, the efficiencies are 97.65\% for the generic-trained tracker and
95.89\% for the baseline. Between 20 and 50~mm, the generic-trained tracker finds 97.10\% of the
targets and the signal-trained one 94.78\%. The gap between the trackers widens with displacement.

\begin{figure*}[tbp]
  \centering
  \fbox{\parbox[c][0.45\linewidth][c]{0.9\textwidth}{\centering placeholder artwork: eff disp}}
  \caption{Efficiency versus the transverse displacement of the production vertex for displaced targets. Solid: the generic-trained tracker; dashed: the signal-trained tracker. The lower panel shows their difference in points.}
  \label{fig:eff_disp}
\end{figure*}

In the bin from 50 to 100~mm, the efficiencies are 95.62\% and 92.62\%. For production radii from
100 to 200~mm, the efficiencies are 94.26\% for the generic-trained tracker and 90.50\% for the
baseline. Between 200 and 300~mm, the generic-trained tracker finds 91.76\% of the targets and the
signal-trained one 87.28\%. In the bin from 300 to 600~mm, the efficiencies are 89.57\% and 83.93\%.
Beyond the outermost pixel layer, both trackers lose particles that leave fewer than five hits. This
loss is a property of the detector, not of the trackers, and it is the same for both. The widening
gap shows that the generic-trained tracker generalizes better precisely where the signal-trained one
should have its advantage.

Figure~\ref{fig:eff_species} breaks the efficiency down by particle species. For muons, the
efficiencies are 97.52\% and 95.61\%, from 2,289 targets. The 1,754 electrons are found with
efficiencies of 97.52\% and 96.16\%. For charged pions, the efficiencies are 96.23\% and 94.76\%,
from 4,843 targets. The 1,533 charged kaons are found with efficiencies of 96.36\% and 94.06\%. For
protons, the efficiencies are 94.04\% and 92.00\%, from 889 targets. Table~\ref{tab:species} lists
the numbers with their bootstrap uncertainties. The lead holds for every species. Electrons are the
hardest case for both trackers, because they lose energy in the material and their tracks bend more
than the network expects.

\begin{figure}[tbp]
  \centering
  \fbox{\parbox[c][0.45\linewidth][c]{0.8\linewidth}{\centering placeholder artwork: eff species}}
  \caption{Efficiency by particle species for the two trackers. The error bars show the bootstrap uncertainty.}
  \label{fig:eff_species}
\end{figure}

\begin{table}[tbp]
  \caption{Efficiency by particle species, with the bootstrap uncertainty in percentage points.}
  \label{tab:species}
  \centering
  \begin{tabular}{lrcc}
    \hline\hline
    Species & Targets & Generic-trained (\%) & Signal-trained (\%) \\
    \hline
    muons & 2,289 & $97.52 \pm 0.31$ & $95.61 \pm 0.27$ \\
    electrons & 1,754 & $97.52 \pm 0.31$ & $96.16 \pm 0.32$ \\
    charged pions & 4,843 & $96.23 \pm 0.27$ & $94.76 \pm 0.27$ \\
    charged kaons & 1,533 & $96.36 \pm 0.37$ & $94.06 \pm 0.26$ \\
    protons & 889 & $94.04 \pm 0.29$ & $92.00 \pm 0.44$ \\
    \hline\hline
  \end{tabular}
\end{table}

Figure~\ref{fig:pileup} shows how the efficiency changes when pile-up is added to the test events.
At an average of 50 pile-up interactions, the generic-trained tracker finds 94.73\% of the targets
and the baseline 91.82\%. With 140 pile-up interactions, the efficiencies fall to 94.17\% and
90.69\%. At 200 interactions per crossing, the efficiencies are 93.72\% and 89.59\%, and the
fake-track rate of the generic-trained tracker rises to 1.12\%. Neither tracker was trained with
pile-up, so both see it here for the first time\footnote{The pile-up events are drawn from the same
generic sample as the training events but from disjoint event numbers.}. The generic-trained tracker
degrades more slowly, which suggests that it has learned the geometry of tracks rather than the
topology of the signal.

\begin{figure}[tbp]
  \centering
  \fbox{\parbox[c][0.45\linewidth][c]{0.8\linewidth}{\centering placeholder artwork: pileup}}
  \caption{Efficiency versus the average number of pile-up interactions. Neither tracker was trained with pile-up.}
  \label{fig:pileup}
\end{figure}

\subsection{Kinematics, signal points and pairs}
\label{sec:kinematics}

The efficiency also depends on the transverse momentum of the particle. Between 1 and 2~GeV, the
generic-trained tracker finds 93.47\% of the targets and the signal-trained one 90.66\%. From 2 to
5~GeV, the efficiencies are 94.28\% and 91.46\%. In the range from 5 to 10~GeV, they are 94.90\% and
92.62\%. Particles with transverse momenta from 10 to 20~GeV are found with efficiencies of 95.13\%
and 92.42\%. For the hardest particles, from 20 to 50~GeV, the efficiencies reach 95.77\% and
93.22\%. Soft particles curl more in the magnetic field, which makes their tracks harder to follow
from layer to layer. The lead of the generic-trained tracker is largest for soft particles, which
are common in generic events and rare in the signal.

The dependence on pseudorapidity is weak for both trackers. Both lose a little efficiency in the
transition between the barrel and the endcaps, where particles cross layers at shallow angles. The
difference between the trackers is the same in the barrel and in the endcaps within its uncertainty.
We therefore do not show it separately.

The headline comparison uses one signal point, but the conclusion should not depend on it. We
therefore repeated the comparison for four combinations of mass and lifetime of the heavy neutral
lepton. For a mass of 10~GeV and a decay length of 100~mm, the efficiencies are 95.41\% and 93.59\%.
At 15~GeV and 100~mm, which is the point the signal-trained tracker was trained on, they are 94.98\%
and 93.17\%. Doubling the decay length to 200~mm at 15~GeV gives 94.64\% and 92.34\%. The heaviest
and longest-lived point, 20~GeV and 300~mm, gives 94.16\% and 91.62\%. The signal-trained tracker
comes closest on its own training point, as expected. Away from it, its efficiency falls faster than
that of the generic-trained tracker.

A search would compare a learned tracker with the standard reconstruction, not only with another
learned tracker. The standard reconstruction with large-radius tracking finds 80.42\% of the
displaced targets at a fake-track rate of 3.36\%. It needs 435~ms per event for the large-radius
pass alone. Both learned trackers find many more displaced targets with far fewer fake tracks. This
comparison is less controlled than the main one, because the standard reconstruction differs in
every stage. We show it only to place the learned trackers on a familiar scale.

Many searches need both decay products of the long-lived particle, because a vertex fit combines
them~\cite{Fischer2022Vertexing}. The generic-trained tracker finds both daughters in 90.27\% of the
signal events, against 84.29\% for the signal-trained tracker. The pair efficiency is roughly the
square of the single-particle efficiency, so small gains per particle become larger gains per event.
The daughters of one decay are often close to each other, and a tracker can merge them into one
track. Merged pairs are rare for both trackers and slightly rarer for the generic-trained one.

\subsection{Failure modes}
\label{sec:failures}

We inspected the displaced targets that one tracker finds and the other misses. Most targets missed
only by the signal-trained tracker have a first hit in the strip layers. In these cases the
signal-trained network assigns low scores to the edges between the first two hits. The
generic-trained network gives the same edges high scores, because it has seen many tracks that start
in the strips from material interactions. Targets missed only by the generic-trained tracker are
fewer and show no common pattern.

Some failures are shared by both trackers. Particles that cross a gap between modules lose a hit and
may be split into two short tracks. Particles that scatter strongly in the material change direction
and are followed only up to the scattering point. Both failure modes depend on the detector and the
graph, not on the training, so they cancel in the comparison. They set the ceiling that neither
tracker can pass without a better graph.

The efficiency averages over particles, but a search selects events. We therefore also counted,
event by event, how many displaced targets each tracker finds. In most signal events both trackers
find the same displaced particles. Where they differ, the generic-trained tracker usually finds one
more particle than the signal-trained tracker. Events in which the signal-trained tracker finds more
are less common and contain mostly soft particles near the edge of the acceptance.

The training size matters for both trackers. With half of the training events, both trackers lose
some efficiency, and the difference between them stays about the same. With a quarter of the
training events, the signal-trained tracker loses more, and the difference grows. The
generic-trained tracker is thus less sensitive to the amount of training data. This is consistent
with the larger number of particles per generic event.

Prompt particles provide a control for the whole comparison. On prompt targets in the signal events,
the two trackers agree to within their uncertainties. The same holds for prompt targets in held-out
generic events. The difference on displaced targets is therefore not a general offset between the
trackers. It appears only for the particles whose production point lies away from the beam line.
This pattern is what one expects if the signal-trained tracker has learned the topology of one decay
rather than the geometry of tracks.

\subsection{Resolution, graph and timing}
\label{sec:resolution}

Figure~\ref{fig:resolution} compares the track parameters with the true values. The transverse
impact parameter is measured with a resolution of 0.018~mm, and the longitudinal one with 0.069~mm.
The relative transverse-momentum resolution is 1.93\%. The azimuthal angle is measured to 0.41~mrad.
These resolutions are the same for both trackers within their uncertainties, because they come from
the hits, not from the network.

\begin{figure}[tbp]
  \centering
  \fbox{\parbox[c][0.45\linewidth][c]{0.8\linewidth}{\centering placeholder artwork: resolution}}
  \caption{Resolution of the track parameters for matched displaced tracks. The two trackers overlap within their uncertainties.}
  \label{fig:resolution}
\end{figure}

Figure~\ref{fig:graph} shows the edge efficiency of the graph as a function of the layer pair. A
perfect edge classifier on the same graph would reach an efficiency of 99.35\%; this is the ceiling
of the pipeline. The generic-trained tracker closes most of the distance between the baseline and
this ceiling. The remaining loss comes from edges that the module map never proposes. Improving the
graph would raise both trackers by the same amount, so it would not change the comparison.

\begin{figure}[tbp]
  \centering
  \fbox{\parbox[c][0.45\linewidth][c]{0.8\linewidth}{\centering placeholder artwork: graph}}
  \caption{Edge efficiency of the constructed graph per pair of adjacent layers. The graph is the same for both trackers.}
  \label{fig:graph}
\end{figure}

Figure~\ref{fig:timing} shows the time spent in each stage of the pipeline on one graphics
processor. Graph construction takes 56~ms per event, the edge classifier 179~ms and track building
24~ms\footnote{The timing excludes reading the events from disk and writing the tracks.}. The total
of 259~ms compares with 296~ms for the standard reconstruction with large-radius tracking enabled.
The two trackers take the same time, because they share the architecture. The timing would allow the
pipeline to run in a software trigger, but we have not optimized it for that purpose.

\begin{figure}[tbp]
  \centering
  \fbox{\parbox[c][0.45\linewidth][c]{0.8\linewidth}{\centering placeholder artwork: timing}}
  \caption{Time per event spent in each stage of the pipeline, on one graphics processor. The standard reconstruction is shown for comparison.}
  \label{fig:timing}
\end{figure}

Figure~\ref{fig:ablation} shows which parts of the network carry the result. Removing the edge
features lowers the efficiency of the generic-trained tracker to 94.78\%. With four message-passing
steps instead of eight, the efficiency drops to 93.96\%. Without the curvature term in the loss, the
efficiency is 94.94\%. The number of message-passing steps matters most, because a displaced track
needs more steps to connect its first hit to the rest. None of the variants reverses the comparison
with the signal-trained tracker.

\begin{figure}[tbp]
  \centering
  \fbox{\parbox[c][0.45\linewidth][c]{0.8\linewidth}{\centering placeholder artwork: ablation}}
  \caption{Efficiency of the generic-trained tracker for variants of the network. The dashed line shows the signal-trained tracker.}
  \label{fig:ablation}
\end{figure}

\subsection{Checks and interpretation}
\label{sec:checks}

We checked that the result does not depend on details of the analysis. Changing the edge threshold
within the range allowed by the validation events moves both efficiencies by the same amount. Using
single instead of double majority raises both efficiencies and leaves the difference unchanged
within its uncertainty. Evaluating on the training process of each tracker shows that both reach a
similar efficiency on their own kind of events. The advantage of the generic-trained tracker
therefore appears only on events that neither tracker was trained on. Each check is repeated for
every training seed, and Appendix~\ref{app:seeds} shows the spread.

Why does the tracker that never saw the signal do better on it? The top-quark pair events contain
many more particles per event than the signal events. The generic-trained network therefore sees
many more examples of tracks, including the few displaced tracks from heavy-flavour decays. The
signal-trained network sees fewer tracks, and most of them share the kinematics of one decay. It
learns that decay well and generalizes less to the different lifetime and mass of the test sample.
This interpretation is consistent with the pile-up study in Fig.~\ref{fig:pileup}, where the
generic-trained tracker degrades more slowly.

Our study has three limitations. The detector is a generic model, and a real detector has dead
modules, misalignment and a more complex material map. The signal is one family of models, and other
long-lived signatures may behave differently. The test events have no pile-up in the headline
comparison, and the pile-up study uses a simplified overlay. None of these limitations favours one
of the trackers, so we expect the comparison to survive them. The numbers in
Table~\ref{tab:headline} and Table~\ref{tab:species} are therefore statements about this simulation
only.

\section{Discussion}
\label{sec:discussion}

The result has a practical consequence for experiments. Signal samples for every model in a search
programme are expensive to simulate, and their number grows with every new idea. Generic samples are
simulated anyway, in large numbers, for every physics analysis. If a tracker trained on generic
samples is at least as good on displaced signals, the expensive signal samples are needed only for
the evaluation. This makes learned tracking easier to adopt and easier to validate.

It also simplifies the statistical treatment of a search. When the tracker is trained on the signal,
the efficiency for a new model must be studied with care, because the tracker may favour the trained
model. A generic-trained tracker has no such preference, so its efficiency for a new model can be
estimated with the same procedure as for any other selection. The uncertainty of that estimate then
comes from the simulation of the detector, not from the training.

We do not claim that generic training is always better. A tracker trained on a large and varied
signal sample might match or surpass the generic-trained one. Our signal sample was as large as the
generic one in events but much smaller in particles. A fair comparison at equal numbers of training
particles is a natural extension of this study. The subset study in the appendix gives a first
indication that the generic events would still win.

Several open questions remain. We have not studied how the result changes with the geometry of the
detector. A detector with fewer layers gives a learned tracker less information per particle and may
change the balance between the trackers. We have also not studied tracks from particles that are not
charged leptons or hadrons, such as long-lived charged particles that decay inside the tracker. Such
particles leave short tracks that end abruptly, which neither training sample contains in large
numbers. Each of these questions can be answered with the same controlled comparison.
